%%=============================================================================
%% Inleiding
%%=============================================================================

\chapter{\IfLanguageName{dutch}{Inleiding}{Introduction}}%
\label{ch:inleiding}

In het moderne voetbal streven professionele clubs continu naar optimalisatie van prestaties, zowel op als naast het veld. Innovaties in trainingsmethoden, medische begeleiding en data-analyse spelen hierbij een steeds grotere rol. Ook binnen scouting wordt in toenemende mate gebruikgemaakt van data om spelers objectief te evalueren en gerichte beslissingen te ondersteunen.

Binnen deze bachelorproef ligt de focus op het scouten van spelers, meer bepaald op de evaluatie van fysieke prestatie-indicatoren zoals afgelegde afstand per 90 minuten, aantal sprints per 90 minuten en andere wedstrijdgerelateerde fysieke statistieken. Binnen de club die meewerkt aan deze bachelorproef leeft de overtuiging dat fysieke factoren een steeds belangrijkere rol spelen in het moderne voetbal. De huidige evaluatie van deze fysieke aspecten gebeurt echter via een doordachte maar deels subjectieve inschatting door de coaching- en scoutingstaff. Hoewel deze aanpak gebaseerd is op expertise, ontbreekt een reproduceerbare en volledig data-gedreven methode om spelers onderling consistent te vergelijken.

Deze subjectiviteit kan leiden tot inconsistenties in beoordeling en maakt het moeilijk om fysieke prestaties objectief te benchmarken over meerdere wedstrijden of tussen verschillende spelers. Er is daarom nood aan een methode die fysieke prestaties op een gestandaardiseerde en reproduceerbare manier kan kwantificeren.

Door middel van machine learning-technieken zal deze bachelorproef onderzoeken hoe fysieke wedstrijddata kan worden gebruikt om een objectieve fysieke spelersscore te ontwikkelen. Hierbij wordt nagegaan welke fysieke prestatie-indicatoren het meest bijdragen aan de constructie van deze score en hoe hun relatieve belang op een interpreteerbare manier kan worden bepaald.

De centrale onderzoeksvraag van deze bachelorproef luidt als volgt: \textbf{Hoe kan een objectieve en reproduceerbare fysieke spelersscore ontwikkeld worden op basis van wedstrijddata met behulp van machine learning-technieken?}

Het einddoel van deze bachelorproef is deze score objectief te kunnen toepassen in de bestaande data-gedreven scoutingstool van de club, zodat deze een waardevolle aanvulling vormt op de huidige evaluatiemethoden en bijdraagt aan het optimaliseren van het scoutingproces.

% Het is gebruikelijk aan het einde van de inleiding een overzicht te
% geven van de opbouw van de rest van de tekst. Deze sectie bevat al een aanzet
% die je kan aanvullen/aanpassen in functie van je eigen tekst.

De rest van deze bachelorproef is als volgt opgebouwd:

In Hoofdstuk~\ref{ch:stand-van-zaken} wordt een overzicht gegeven van de stand van zaken binnen het onderzoeksdomein, op basis van een literatuurstudie.

In Hoofdstuk~\ref{ch:methodologie} wordt de methodologie toegelicht en worden de gebruikte onderzoekstechnieken besproken om een antwoord te kunnen formuleren op de onderzoeksvragen.

% TODO: Vul hier aan voor je eigen hoofstukken, één of twee zinnen per hoofdstuk

In Hoofdstuk~\ref{ch:conclusie}, tenslotte, wordt de conclusie gegeven en een antwoord geformuleerd op de onderzoeksvragen. Daarbij wordt ook een aanzet gegeven voor toekomstig onderzoek binnen dit domein.